%!TEX program = xelatex
% ============================================================
% T01 - Metropolis · Ocean Gradient
% 风格：深蓝 + 青色，科技感，现代扁平
% 适合：科技产品介绍、商务提案、项目汇报
% 验证：Overleaf XeLaTeX ✅ (2026-03-05)
% ============================================================

\documentclass[aspectratio=169, 10pt]{beamer}

\usetheme{metropolis}

% ---------- 中文支持 ----------
\usepackage[UTF8]{ctex}
\usepackage{fontspec}

% ---------- 常用宏包 ----------
\usepackage{booktabs}
\usepackage{tabularx}
\usepackage{xcolor}
\usepackage{array}

\setmonofont{Courier New}

% ---------- 配色：Ocean Gradient ----------
\definecolor{navyBlue}{HTML}{1E2761}
\definecolor{deepBlue}{HTML}{065A82}
\definecolor{tealAccent}{HTML}{14B8A6}
\definecolor{offWhite}{HTML}{F0F4F8}
\definecolor{mutedText}{HTML}{64748B}
\definecolor{goldText}{HTML}{F59E0B}

\setbeamercolor{normal text}{fg=navyBlue, bg=offWhite}
\setbeamercolor{alerted text}{fg=tealAccent}
\setbeamercolor{example text}{fg=deepBlue}
\setbeamercolor{frametitle}{bg=navyBlue, fg=white}
\setbeamercolor{progress bar}{fg=tealAccent, bg=navyBlue!30}
\setbeamercolor{title separator}{fg=tealAccent}
\setbeamercolor{structure}{fg=navyBlue}
\setbeamercolor{block title}{bg=navyBlue, fg=white}
\setbeamercolor{block body}{bg=offWhite}
\setbeamercolor{block title alerted}{bg=tealAccent, fg=white}
\setbeamercolor{block title example}{bg=deepBlue, fg=white}

\metroset{
  progressbar       = frametitle,
  sectionpage       = progressbar,
  block             = fill,
}

% ---------- 页脚 ----------
\setbeamertemplate{footline}{
  \begin{beamercolorbox}[wd=\paperwidth, ht=2.5ex, dp=1ex]{frametitle}
    \hfill
    \color{white!70}\footnotesize
    % ← 在此填写机构名
    \quad|\quad
    \insertframenumber\,/\,\inserttotalframenumber
    \hspace{1em}
  \end{beamercolorbox}
}

% ============================================================
% 填写元信息
% ============================================================
\title{[演讲标题]}
\subtitle{[副标题]}
\author{[作者]}
\institute{[机构]}
\date{[日期]}

% ============================================================
\begin{document}
% ============================================================

% ---- 深色封面 ----
{
  \setbeamercolor{background canvas}{bg=navyBlue}
  \begin{frame}
    \titlepage
    \vspace{-0.5em}
    \begin{center}
      \textcolor{tealAccent}{\small [封面副文字]}
    \end{center}
  \end{frame}
}

% ---- 目录（可选，短演讲可删） ----
\begin{frame}{目录}
  \setbeamertemplate{section in toc}[sections numbered]
  \tableofcontents[hideallsubsections]
\end{frame}

% ============================================================
\section{[第一节]}
% ============================================================

% --- 三栏卡片布局 ---
\begin{frame}{[标题]}
  \begin{columns}[T]
    \begin{column}{0.31\textwidth}
      \begin{block}{[卡片1标题]}
        \small [内容]
      \end{block}
    \end{column}
    \begin{column}{0.31\textwidth}
      \begin{alertblock}{[卡片2标题]}
        \small [内容]
      \end{alertblock}
    \end{column}
    \begin{column}{0.31\textwidth}
      \begin{exampleblock}{[卡片3标题]}
        \small [内容]
      \end{exampleblock}
    \end{column}
  \end{columns}
\end{frame}

% --- 表格布局 ---
\begin{frame}{[标题]}
  \begin{tabularx}{\textwidth}{lX}
    \toprule
    \textbf{列A} & \textbf{列B} \\
    \midrule
    行1 & 内容 \\
    行2 & 内容 \\
    \bottomrule
  \end{tabularx}
\end{frame}

% --- 两栏布局 ---
\begin{frame}{[标题]}
  \begin{columns}[c]
    \begin{column}{0.45\textwidth}
      [左侧内容]
    \end{column}
    \begin{column}{0.50\textwidth}
      [右侧内容]
    \end{column}
  \end{columns}
\end{frame}

% --- 层叠块（架构/流程）---
\begin{frame}{[标题]}
  \begin{block}{[顶层]}\small [内容]\end{block}
  \vspace{0.3em}
  \begin{alertblock}{[中层]}\small [内容]\end{alertblock}
  \vspace{0.3em}
  \begin{exampleblock}{[底层]}\small [内容]\end{exampleblock}
\end{frame}

% ============================================================
\section{[第二节]}
% ============================================================

\begin{frame}{[标题]}
  \begin{itemize}
    \item [要点1]
    \item [要点2]
    \item [要点3]
  \end{itemize}
\end{frame}

% ---- 深色结束页 ----
{
  \setbeamercolor{background canvas}{bg=navyBlue}
  \begin{frame}[standout]
    \textcolor{tealAccent}{\Large 谢谢}\\[0.5em]
    \textcolor{white}{\small [机构名]}\\
    \textcolor{white!60}{\footnotesize [日期]}
  \end{frame}
}

\end{document}
